\section{Constructive kinetic bounds and their precise
limits}\label{constructive-kinetic-bounds-and-their-precise-limits}

This section proves a second-order correction to the full-orbit
weak-quasisymmetry momentum and a finite-time, pointwise hot-core
theorem for linear Vlasov transport in a prescribed field with an
absorbing loss boundary. The theorem retains Maxwellian velocity tails
and permits a quantitative perturbation of the initial distribution. It
is not a result for self-consistent Vlasov--Maxwell--Landau. In
particular, it does not repair the equilibrium obstruction of the
original weak-QS seed. The new arbitrary harmonic vacuum-field
construction does not automatically satisfy weak QS and cannot be
substituted into the hypotheses without a separate proof.

\subsection{A corrected flux for full Lorentz
orbits}\label{a-corrected-flux-for-full-lorentz-orbits}

\paragraph{Hypotheses and
conventions}\label{hypotheses-and-conventions}

Let \(\Omega\subset\mathbb R^3\) be a bounded domain with \(C^2\)
boundary. Suppose \(B,u,\psi\) are \(C^3\) on a neighborhood of
\(\overline\Omega\), \(B\) never vanishes, and

\[
\nabla\cdot B=\nabla\cdot u=0,\qquad
B\times u=\nabla\psi,\qquad u\cdot\nabla |B|^2=0.
\tag{Q1}
\]

Write \(\beta=|B|,\ b=B/\beta,\ P=I-bb^T\), and

\[
S=\operatorname{sym}\nabla u,\qquad
K v=v\times b.
\]

Thus \(K^T=-K\) and \(K^2=-P\). All tensors and norms below are
Euclidean tensors and norms in physical space, not coordinate-component
norms in the embedded-torus chart.

Taking curl in \(B\times u=\nabla\psi\), and using both divergence
conditions, gives \([B,u]=0\). Consequently

\[
\operatorname{tr}S=0,\qquad
b^TSb=\frac{B\cdot((B\cdot\nabla)u)}{\beta^2}
     =\frac{u\cdot\nabla\beta^2}{2\beta^2}=0.
\tag{Q2}
\]

No assumptions about closed \(u\)-orbits, dense \(B\)-orbits, or
regularity of every interior flux level are required for the following
algebra.

\paragraph{Solving the gyrophase homological equation
explicitly}\label{solving-the-gyrophase-homological-equation-explicitly}

Put

\[
a=Sb,\qquad S_\perp=PSP,\qquad
C=b(Ka)^T+(Ka)b^T-\frac12S_\perp K,\qquad D=C/\beta.
\tag{Q3}
\]

Then \(a\perp b\), \(S_\perp\) has trace zero on \(b^\perp\), and

\[
C=C^T,\qquad [C,K]=CK-KC=S.
\tag{Q4}
\]

\textbf{Proof.} Decompose

\[
S=S_\perp+ba^T+ab^T.
\]

On the two-dimensional plane \(b^\perp\), a symmetric trace-zero matrix
anticommutes with the quarter-turn \(K\): \(S_\perp K+KS_\perp=0\). Thus
\(S_\perp K\) is symmetric. If \(C_1=b(Ka)^T+(Ka)b^T\), direct
multiplication using \(Kb=0\) and \(K^2a=-a\) gives
\([C_1,K]=ba^T+ab^T\). The same identities give
\([-\frac12S_\perp K,K]=S_\perp\). This proves both claims. In Frobenius
norm,

\[
\|C\|_F^2=2|a|^2+\frac14\|S_\perp\|_F^2
\le \|S\|_F^2.
\]

For any symmetric \(C\),

\[
\mathcal R_b(v^TCv)
:=(v\times b)\cdot\nabla_v(v^TCv)
=v^T(CK-KC)v.
\]

Therefore (Q4) is exactly the homological equation
\(\mathcal R_b(v^TCv)=v^TSv\). It is globally tensorial and requires no
choice of a gyroangle or frame. \(\square\)

\paragraph{Exact full-orbit
identity}\label{exact-full-orbit-identity}

For a species of mass \(m>0\) and charge \(q\ne0\), write the signed
quantity \(\varepsilon=m/q\). Define

\[
J(x,v)=\psi(x)-\varepsilon u(x)\cdot v
               +\varepsilon^2v^TD(x)v.
\tag{Q5}
\]

Consider prescribed-field Lorentz trajectories

\[
\dot x=v,\qquad
\dot v=\varepsilon^{-1}v\times B(x).
\tag{Q6}
\]

Then, exactly,

\[
\frac{d}{dt}J(x(t),v(t))
 =\varepsilon^2
   \sum_{i,j,k}v_i v_j v_k\,\partial_{x_k}D_{ij}(x(t)).
\tag{Q7}
\]

\textbf{Proof.} The first two terms of (Q5) give

\[
\frac d{dt}(\psi-\varepsilon u\cdot v)
 =v\cdot\nabla\psi-\varepsilon v^TSv-u\cdot(v\times B)
 =-\varepsilon v^TSv.
\]

For the last term,

\[
\frac d{dt}\big(\varepsilon^2v^TDv\big)
 =\varepsilon^2\sum_{ijk}v_iv_jv_k\partial_kD_{ij}
  +\varepsilon\beta\,\mathcal R_b(v^TDv).
\]

Since \(D=C/\beta\), the last summand equals \(\varepsilon v^TSv\) by
(Q4), and cancels the preceding defect. \(\square\)

This improves the exact first-order weak-QS momentum defect by one
order. It does not make \(J\) an exact invariant. The remainder is
explicitly computable from the physical-space derivatives of the seed.

\paragraph{A deterministic confinement
certificate}\label{a-deterministic-confinement-certificate}

Define

\[
U_0=\sup_{\overline\Omega}|u|,\qquad
D_0=\sup_{\overline\Omega}\|D\|_{\rm op},\qquad
D_1=\sup_{\substack{x\in\overline\Omega\\ |w|=1}}
\left|\sum_{ijk}w_iw_jw_k\partial_kD_{ij}(x)\right|.
\]

For \(V>0\), set

\[
A(V)=|\varepsilon|U_0V+\varepsilon^2D_0V^2,\qquad
R(V)=\varepsilon^2D_1V^3.
\]

For any trajectory segment remaining in \(\overline\Omega\), with
\(|v|=V_*\le V\),

\[
|\psi(x(t))-\psi(x(s))|
 \le 2A(V)+|t-s|R(V).
\tag{Q8}
\]

Indeed, \(|J-\psi|\le A(V)\), and integrating (Q7) gives
\(|J(t)-J(s)|\le |t-s|R(V)\). The speed is constant by (Q6). The same
argument works forwards and backwards in time.

Suppose in addition \(\psi=\psi_b\) on \(\partial\Omega\),
\(\psi<\psi_b\) in \(\Omega\), and define a core

\[
\Omega_\Delta=\{x\in\Omega:\psi(x)\le\psi_b-\Delta\}.
\]

Interior critical points of \(\psi\), including a magnetic axis, are
allowed. If

\[
\Delta>2A(V)+TR(V),
\tag{Q9}
\]

then every characteristic through \(x\in\Omega_\Delta,\ |v|\le V\), at
any time \(0\le t\le T\), remains inside \(\Omega\) when traced back to
time zero. For otherwise stop at the first backward boundary hit and
apply (Q8) on the segment before that hit. The change in \(\psi\) would
be at least \(\Delta\), contradicting (Q9). The strict inequality also
excludes a grazing boundary hit. The analogous forward no-hit assertion
holds for initial points in the core.

If \(R(V)>0\), this gives the explicit time

\[
T_{\rm cert}=\frac{\Delta-2A(V)}{R(V)}
\]

whenever its numerator is positive. If the numerator is nonpositive,
this certificate gives no information; its failure is not an escape
theorem. If \(R(V)=0\) and the numerator is positive, the algebraic
certificate holds for arbitrary time in the prescribed-field model.

\paragraph{Scaling}\label{scaling}

Under \(B\mapsto\Lambda B,\ \psi\mapsto\Lambda\psi,\ u\mapsto u\), one
has \(C\mapsto C,\ D\mapsto D/\Lambda\), while the geometric flux buffer
scales by \(\Lambda\). Consequently the cubic accumulated drift relative
to the flux buffer is of order \(T/\Lambda^2\). This is a precise
strong-field gain in this orbit certificate, not a statement about
self-consistent plasma stability at arbitrary field amplitude. No
numerical pass at the project's 10 T budget has been established here:
\(U_0,D_0,D_1\), and the physical flux buffer must first be bounded for
the actual embedded geometry.

\paragraph{Exact relativistic orbit
extension}\label{exact-relativistic-orbit-extension}

The orbit identity also has a direct relativistic version. Let \(w=p/m\), let
\(c_{\rm light}\) denote the speed of light, and set
\(\gamma(w)=(1+|w|^2/c_{\rm light}^2)^{1/2}\). For prescribed \(E=0\),
\[
\dot x=w/\gamma,\qquad
\dot w=(\varepsilon\gamma)^{-1}w\times B.
\] Both \(|w|\) and \(\gamma\) are constant. The vector field on
\((x,w)\) is exactly \(1/\gamma(w)\) times the Newtonian vector field
used in (Q6). Hence \[
J_{\rm rel}=\psi-\varepsilon u\cdot w+\varepsilon^2w^TDw,\qquad
\dot J_{\rm rel}=\frac{\varepsilon^2}{\gamma}
\sum_{ijk}w_iw_jw_k\partial_kD_{ij}.
\] On \(|w|\le W\), the endpoint bound is \(A(W)\), and the drift bound
is \[
R_{\rm rel}(W)=
\frac{\varepsilon^2D_1W^3}
{\sqrt{1+W^2/c_{\rm light}^2}}.
\] The function \(r^3/\sqrt{1+r^2/c_{\rm light}^2}\) is increasing for
\(r\ge0\), which justifies this uniform bound. The no-hit proof is
unchanged. The Maxwellian temperature theorem below remains explicitly
Newtonian; no identification of its scalar moment temperature with a
relativistic Jüttner rest-frame temperature is made.

\subsection{A pointwise hot-core theorem at an absorbing
boundary}\label{a-pointwise-hot-core-theorem-at-an-absorbing-boundary}

Consider the prescribed-field Vlasov equation

\[
\partial_t f+v\cdot\nabla_xf+
\varepsilon^{-1}(v\times B)\cdot\nabla_vf=0
\]

on \(\Omega\times\mathbb R^3\), with zero incoming data at the loss
boundary. Take

\[
f(0,x,v)=M(v)
=n_0(2\pi\sigma^2)^{-3/2}e^{-|v|^2/(2\sigma^2)},
\qquad \sigma^2=\Theta_0/m.
\tag{H1}
\]

Here \(\Theta_0\) is temperature in energy units. The initial trace and
zero incoming boundary data are not compatible at the time-zero corner;
the assertion is about the standard mild solution defined by backward
characteristics killed at their first boundary encounter, not a smooth
solution up to that corner.

Along a characteristic that reaches the initial surface, the value is
\(M(v(0))=M(v)\), since speed is conserved. Along one killed at the
boundary, the value is zero. Consequently

\[
0\le f(t,x,v)\le M(v)
\quad\hbox{everywhere},
\]

and (Q9) gives the stronger equality

\[
f(t,x,v)=M(v)
\quad\hbox{for }x\in\Omega_\Delta,\ |v|\le V,\ 0\le t\le T.
\tag{H2}
\]

No compact velocity support has been assumed. High-velocity particles
can be lost.

Define the local density, mean velocity and scalar kinetic temperature
by

\[
n=\int f\,dv,\quad U=n^{-1}\int vf\,dv,\quad
\Theta=\frac{m}{3n}\int |v-U|^2f\,dv.
\]

Let \(c=V/\sigma\), and define

\[
a(c)=\operatorname{erf}(c/\sqrt2)
       -\sqrt{2/\pi}\,c e^{-c^2/2},
\] \[
b(c)=\operatorname{erf}(c/\sqrt2)
       -\sqrt{2/\pi}\,(c+c^3/3)e^{-c^2/2},
\] \[
d(c)=\sqrt{2/\pi}(c^2+2)e^{-c^2/2}.
\]

Equivalently \(a=\mathbb P(\chi_3\le c)\) and
\(b=\mathbb P(\chi_5\le c)\). On the core and certified time interval,

\[
a(c)n_0\le n\le n_0,\qquad
\boxed{\quad
\frac{\Theta}{\Theta_0}
\ge b(c)-\frac{d(c)^2}{3a(c)^2}.
\quad}
\tag{H3}
\]

\textbf{Proof.} Integrating (H2) proves the lower density bound and
\(\int |v|^2f\ge3n_0\sigma^2b\); the global upper bound gives
\(n\le n_0\). The truncated Maxwellian has zero first moment. Thus the
entire mean-momentum contribution comes from the remaining tail, and

\[
\left|\int vf\,dv\right|
\le\int_{|v|>V}|v|M\,dv=n_0\sigma d.
\]

Using \(\Theta=m(\int |v|^2f/n-|\int vf|^2/n^2)/3\) and the density
bounds proves (H3). This explicitly allows the lost tail to be
asymmetric, and therefore does not assume zero mean flow after losses.
\(\square\)

For \(c=2.8\), the numerical constants are

\[
a>0.95056321,\quad
b-\frac{d^2}{3a^2}>0.82577125.
\]

Hence \(\Theta_0=12.5\,{\rm keV}\) implies \(\Theta>10.3221\,{\rm keV}\)
and \(n>0.95056n_0\) on the certified core and interval. For \(c=3\),
the temperature ratio is greater than \(0.88757292\). These strict
rational lower bounds are independently checked by the accompanying
alternating-series certificate. These are
implications of (Q9), not a claim that (Q9) has passed for the actual
device dimensions.

\paragraph{Stability under initial distribution
perturbations}\label{stability-under-initial-distribution-perturbations}

Suppose instead \[
(1-\delta)M\le f_0(x,v)\le(1+\delta)M,\qquad 0\le\delta<1.
\]

The same characteristic argument gives the same bounds for \(|v|\le V\)
in the core, and \(0\le f\le(1+\delta)M\) globally. If
\(d_0=2\sqrt{2/\pi}\), then

\[
n\ge(1-\delta)a n_0,\qquad
\frac{\Theta}{\Theta_0}\ge
\frac{1-\delta}{1+\delta}b
-\frac{(d+\delta d_0)^2}{3(1-\delta)^2a^2}.
\tag{H4}
\]

To prove the momentum bound used here, subtract the Maxwellian on the
protected ball. Its perturbation contributes at most
\(\delta n_0\sigma(d_0-d)\); the unprotected tail contributes at most
\((1+\delta)n_0\sigma d\), whose sum is \(n_0\sigma(d+\delta d_0)\). The
second-moment and density estimates are as before. At
\(c=2.8,\delta=0.01\), (H4) gives
\(\Theta/\Theta_0>0.80709315\), hence \(10.08866\,{\rm keV}\)
from \(12.5\,{\rm keV}\).

There is also a directional temperature bound, so this conclusion need
not hide a cold direction behind a large scalar trace. Let \[
\mathsf T=\frac mn\int(v-U)(v-U)^Tf\,dv.
\] The protected Maxwellian ball has second-moment matrix
\(n_0\sigma^2bI\). The same estimates therefore give \[
\frac{\mathsf T}{\Theta_0}\ \ge\
\left[\frac{1-\delta}{1+\delta}b
-\frac{(d+\delta d_0)^2}{(1-\delta)^2a^2}\right]I
\tag{H5}
\] as a quadratic-form inequality. For \(c=2.8,\delta=0\), its
coefficient is \(0.8078673232710732\), giving every directional
temperature above \(10.0983\,{\rm keV}\) from \(12.5\,{\rm keV}\). For
\(c=3,\delta=0.01\), its coefficient is \(0.8593547807238485\), giving
\(10.7419\,{\rm keV}\).

This is finite-time stability for an explicitly stated transport
problem. Electromagnetic feedback, Coulomb collisions, fusion, alpha
heating and fueling are absent from its hypotheses.

\subsection{Electric fields, collisions, and sources: the exact
additional
terms}\label{electric-fields-collisions-and-sources-the-exact-additional-terms}

Keep \(J\) constructed from the static reference weak-QS field \(B\).
Let the actual Lorentz field be \(B+\delta B(t,x)\), and include
\(E(t,x)\). Then

\[
\frac{dJ}{dt}
=\varepsilon^2\sum v_iv_jv_k\partial_kD_{ij}
 +(-u+2\varepsilon Dv)\cdot(E+v\times\delta B).
\tag{P1}
\]

This is an exact identity. In particular an electrostatic field constant
on flux surfaces need not annihilate the second-order term, even though
\(u\cdot E=0\).

The speed obeys \[
|v(s)|\le |v(t)|+|\varepsilon|^{-1}
\int_s^t\|E(r)\|_\infty\,dr
\] backwards (and the corresponding forward bound). For a known bound
\(V_*\) on the relevant speed range,

\[
|\dot J|\le
\varepsilon^2D_1V_*^3+
(U_0+2|\varepsilon|D_0V_*)
\big(\|E\|_\infty+V_*\|\delta B\|_\infty\big).
\tag{P2}
\]

Thus the orbit certificate survives with the integral of the right side
and \(A(V_*)\). A self-consistent theorem still has to prove these field
budgets from the initial data and actual sources; assuming them for all
time would assume the missing part of stability. With \(E\ne0\), the
Maxwellian bounds (H2)--(H4) also require accounting for the
initial-to-final speed change; they do not follow unchanged merely from
(P2).

For a prescribed kinetic collision generator in backward-observable form
\[
\mathcal L_v\phi=\ell(x,v)\cdot\nabla_v\phi
                 +a(x,v):D_v^2\phi,\qquad a\ge0,
\] one has \[
\mathcal L_vJ
=\ell\cdot(-\varepsilon u+2\varepsilon^2Dv)
 +2\varepsilon^2a:D.
\tag{P3}
\]

In an equivalent diffusion representation with covariance \(2a\), the
martingale part of \(J\) has quadratic-variation rate \[
2(-\varepsilon u+2\varepsilon^2Dv)^T
 a(-\varepsilon u+2\varepsilon^2Dv).
\tag{P4}
\]

Therefore collisions do not preserve the corrected invariant or the
deterministic protected-velocity ball. A stochastic first-exit estimate
needs the actual drift, diffusion, and high-speed tail bounds. The full
nonlinear Landau coefficients depend on the unknown distributions;
replacing them by fixed Maxwellian coefficients is a different model.

For completeness, there is a useful exact weak-form Landau estimate. For
one-species Coulomb Landau with kernel
\(a_C(w)=\kappa|w|^{-1}(I-\hat w\hat w^T)\), set \[
\mathcal D(f)=\frac12\iint ff'
(\nabla\log f-\nabla'\log f')^Ta_C(v-v')
(\nabla\log f-\nabla'\log f')\,dv\,dv'.
\] For a \(C^2\) observable \(\phi\) with
\(\|D_v^2\phi\|_{\rm op}\le H\), \[
\left|\int Q(f,f)\phi\,dv\right|
\le \sqrt{\kappa}\,H\,
n^{3/4}\left(\int|v|^2f\,dv\right)^{1/4}
\sqrt{\mathcal D(f)}.
\tag{P5}
\]

Indeed the weak form is minus one half of the pair integral of
\(ff'(\nabla\log f-\nabla'\log f')^Ta_C(\nabla\phi-\nabla'\phi')\).
Cauchy--Schwarz gives the square root of \(\mathcal D\) times the square
root of \(\frac12\iint ff'Y^Ta_CY\). The Lipschitz bound
\(|Y|\le H|v-v'|\) reduces the latter to
\((\kappa H^2/2)\iint ff'|v-v'|\), and
\(\iint ff'|v-v'|\le2n^{3/2}(\int|v|^2f)^{1/2}\).

For \(\phi=J\), the Hessian is \(2\varepsilon^2D\), so the collision
effect on the mean corrected flux is genuinely second order. First-order
collision contributions vanish because self-collisions conserve local
momentum. This is a mean-observable estimate, not an exit-probability or
temperature theorem. Spatial integration of (P5) still needs a density
bound (for example \(n\le N\)), or a substitute estimate. With that
bound, total mass \(M\), and second moment \(E_2\), \[
\left|\iint Q(f,f)\phi\,dv\,dx\right|
\le \sqrt{\kappa}\,H\sqrt N\,(ME_2)^{1/4}
       \sqrt{\int_\Omega\mathcal D(f)\,dx}.
\] It is essential not to declare \(N\) controlled for the nonlinear
open-boundary system without proving it. For multispecies collisions the
individual species momentum is not conserved; the mass-weighted sum is,
and the species-specific corrected fluxes do not automatically match
those weights.

For an added distribution source \(S_f\), the exact observable balance
acquires \(\iint J S_f\); outgoing absorbing particles also carry the
flux \(Jf(v\cdot n)\). Scalar heating power does not specify this source
moment or its entropy cost.

\subsection{Covariant derivatives and the remaining magnetic
shear}\label{covariant-derivatives-and-the-remaining-magnetic-shear}

For \(\mathcal R_a=(v\times a)\cdot\nabla_v\), one has \[
[\mathcal R_a,\mathcal R_c]=\mathcal R_{a\times c}.
\] Define \[
a_i=(\partial_i b)\times b,\qquad
\mathfrak D_i=\partial_{x_i}+\mathcal R_{a_i},\qquad
\Omega_c=q\beta/m.
\] The vector identity \(a_i\times b=-\partial_i b\) gives \[
[\mathfrak D_i,\Omega_c\mathcal R_b]
 =(\partial_i\Omega_c)\mathcal R_b.
\tag{C1}
\] Thus the derivative of the \textbf{direction} of the magnetic field
cancels exactly. This removes its factor \(\beta\). It leaves the
derivative of the cyclotron frequency.

The remaining exact transport commutators are \[
[\mathfrak D_i,v\cdot\nabla_x]
 =(v\times a_i)\cdot\nabla_x
   -\mathcal R_{(v\cdot\nabla_x)a_i},
\tag{C2}
\] \[
[\mathfrak D_i,(q/m)E\cdot\nabla_v]
 =(q/m)(\partial_iE+a_i\times E)\cdot\nabla_v.
\tag{C3}
\] In (C2), the last rotation parameter is itself linear in \(v\), so
its coefficient before \(\nabla_v\) is quadratic in velocity. The
reduction of a large field-direction coefficient is useful, but it does
not eliminate geometric derivatives or velocity-weight requirements.

For this particular minimal-rotation connection, \[
[\mathfrak D_i,\mathfrak D_j]
 =-[\,b\cdot(\partial_i b\times\partial_j b)\,]\mathcal R_b.
\tag{C4}
\] To check (C4), differentiate \(a_i=\partial_i b\times b\); the
derivative part is \(-2\partial_i b\times\partial_j b\), and
\(a_i\times a_j=\partial_i b\times\partial_j b\). The cross product is
parallel to \(b\). One can modify the connection by adding
\(\alpha_i b\). A full chosen orthonormal frame can give a flat
connection locally; thus (C4) should not be advertised as an
unconditional topological obstruction. The corresponding streaming
connection terms remain, and (C1)'s magnitude derivative is unaffected
by such a frame choice.

Rotational covariance of the Coulomb kernel gives \[
\mathfrak D_iQ(f,g)=Q(\mathfrak D_i f,g)+Q(f,\mathfrak D_i g).
\tag{C5}
\] For a spatially uniform isotropic Maxwellian reference, the
linearized Landau operator therefore commutes with \(\mathfrak D_i\). It
does not follow that a linearization around a spatially varying
anisotropic reference commutes.

The remaining magnitude shear cannot be removed by a smooth periodic
gyroangle gauge. At fixed speed and pitch, write the rotation as
\(\Omega_c(x)\partial_\vartheta\). A derivative
\(\partial_i+\alpha(x,\vartheta)\partial_\vartheta\) would commute with
it only if \[
\partial_\vartheta\alpha=\partial_i\log|\Omega_c|.
\] Integrating around the gyroangle circle forces
\(\partial_i\Omega_c=0\). The pure rotation subproblem already exhibits
the shear: \[
g(t,x,\vartheta)=g_0(x,\vartheta-\Omega_c(x)t)
\] has a spatial derivative containing
\(-t(\partial_i\Omega_c)\partial_\vartheta g_0\). This is an obstruction
to that proposed cancellation, not a theorem that the full collisional
transport system is unstable. A time-dependent transported derivative
can cancel the shear, but its coefficients grow in time and a uniformly
equivalent norm still needs proof.

\paragraph{A global rotating frame on the solid
torus}\label{a-global-rotating-frame-on-the-solid-torus}

On a solid torus the oriented two-plane bundle \(b^\perp\) is trivial:
oriented two-plane bundles are classified by their Euler class in
\(H^2(\Omega;\mathbb Z)=0\). Thus a smooth global \(R(x)\in SO(3)\) with
\(R(x)e_3=b(x)\) exists. Put \[
A_i=R^T\partial_iR,\qquad g(x,w)=f(x,Rw).
\] Each \(A_i\) is skew. The prescribed transport becomes exactly \[
\partial_tg+(Rw)\cdot\nabla_xg+
\left[\Omega_c(w\times e_3)
      -\sum_i(Rw)_iA_iw\right]\cdot\nabla_wg.
\tag{C6}
\] This follows by the chain rule:
\(\partial_i f=\partial_i g-(A_iw)\cdot\nabla_wg\) at fixed physical
velocity. The physical covariant derivatives corresponding to fixed
\(w\) form a flat pure-gauge connection. Their rotation parameters equal
\(\partial_i b\times b+\alpha_i b\) for suitable \(\alpha_i\), so they
still commute with \(\mathcal R_b\).

The streaming and connection terms in (C6) satisfy the exact divergence
cancellation \[
\nabla_x\cdot(Rw)
+\nabla_w\cdot\left(-\sum_i(Rw)_iA_iw\right)=0.
\tag{C7}
\] Indeed the first term is \(\sum_i(R^Te_i)^TA_iw\), and the second is
its negative because \(\operatorname{tr}A_i=0\). Also \(w\cdot A_iw=0\).
Every radial velocity weight is therefore invariant under the connection
transport, and this transport is skew in weighted \(L^2(dx\,dw)\), up to
the actual spatial boundary flux. The Jacobian of
\((x,w)\mapsto(x,R(x)w)\) is one. Landau's rotational covariance
preserves its collision form.

Under \(B\mapsto\Lambda B\), \(R,A_i\) remain unchanged. Only
\(\Omega_c\) is multiplied by \(\Lambda\). Differentiation still
produces \((\partial_i\Omega_c)\mathcal R_{e_3}g\), derivatives of \(R\)
multiplying \(|w|\nabla_xg\), and derivatives of \(A_i\) multiplying
\(|w|^2\nabla_wg\). This global-frame formulation removes the
field-direction large parameter at zero order and displays where
weighted high-order estimates must still close.

\paragraph{A non-tautological weighted estimate for a simpler
collision
operator}\label{a-non-tautological-weighted-estimate-for-a-simpler-collision-operator}

There is a useful constructive check that the remaining frequency shear
need not imply instability. Freeze spatial transport and use the
normalized Ornstein--Uhlenbeck generator \(L=\Delta_w-w\cdot\nabla_w\)
on Gaussian \(L^2(M\,dw)\): \[
\partial_tg+\Omega_c(x)\mathcal R_{e_3}g=\nu Lg.
\] For one spatial derivative put \(h=\partial_i g\),
\(a(x)=1+|\partial_i\Omega_c|^2/\nu^2\), and \[
\mathcal E_i=\int\big(|h|^2+a|\mathcal R_{e_3}g|^2\big)M\,dw\,dx.
\] Commutation of \(L\) with rotations, Gaussian integration by parts,
and Gaussian Poincare give \[
\frac d{dt}\mathcal E_i
\le-\nu\int|\nabla_wh|^2M
    -\nu\int(a+1)|\nabla_w\mathcal R_{e_3}g|^2M.
\tag{C8}
\] To verify this, the only cross term is
\(-2(\partial_i\Omega_c)\langle h,\mathcal R_{e_3}g\rangle_M\). The
rotation has zero Gaussian mean. Subtract the Gaussian mean of \(h\),
apply Poincare and Young, and then apply Poincare to
\(\mathcal R_{e_3}g\). The cross term costs at most
\(\nu\|\nabla_wh\|^2+\nu(a-1)\|\nabla_w\mathcal R_{e_3}g\|^2\), which
yields (C8).

This is an actual modified energy estimate, but only for the stated
subproblem. Restoring spatial streaming differentiates \(a(x)\) and
introduces the connection terms in (C6). Full Landau has momentum and
energy null modes and weak high-velocity coercivity, so the two Gaussian
Poincare substitutions do not give (C8) for that operator. Those are
concrete missing estimates, not an inference that no such estimate can
exist.

For a vacuum field, \(\nabla\cdot B=\nabla\times B=0\) implies
\(\Delta B=0\), hence \[
\Delta |B|^2=2|\nabla B|^2.
\] Consequently a vacuum field with constant magnitude is constant as a
vector field. A nontrivial vacuum stellarator cannot generally evade
(C1) by choosing a nonconstant field of constant magnitude.

\subsection{What remains unproved}\label{what-remains-unproved}

The finite-time transport theorem is a genuine consequence of explicit
geometric bounds, provided its buffer inequality is satisfied. We have
not certified that inequality for the project's actual embedded geometry
and field budget. The common-temperature Maxwellian used here has no
self-consistent field response, and the prescribed original weak-QS
field has a separately established obstruction to the proposed
equilibrium classes. Thus this construction is an orbit-level and
kinetic-transport advance, not an externally coiled hot fusion
equilibrium. Extending it requires actual collision and source
estimates, electromagnetic feedback bounds, wall physics and a globally
coercive continuation estimate. The covariant calculation removes the
large field-direction commutator but leaves field-strength shear,
spatial transport terms and velocity-weight requirements.

\subsection{Primary sources and method
provenance}\label{primary-sources-and-method-provenance}

The explicit corrected flux and absorbing-core inequalities above are calculations made here. The following primary references delimit adjacent established results; the geometric inspiration is discussed earlier in this note.

Karaki's arbitrary-smooth-field hypocoercivity result concerns a linear
Fokker--Planck operator on the flat torus:

\url{https://doi.org/10.3934/krm.2020011}

Fan--Lei's full Vlasov--Maxwell--Landau result allows a strong uniform
background magnetic field:

\url{https://doi.org/10.1007/s10955-022-03030-1}

Dong--Guo--Yastrzhembskiy's global bounded-domain relativistic
Vlasov--Maxwell--Landau theorem is near Jüttner, with electromagnetic
smallness and zero magnetic flux through toroidal cuts:

\url{https://arxiv.org/abs/2401.00554}

Ouyang--Wu's October 2026 preprint gives a global near-Maxwellian
theorem for the Landau equation with inflow or diffuse reflection, but
does not assert a full strong-variable-field Vlasov--Maxwell--Landau
theorem:

\url{https://arxiv.org/abs/2610.40517}

Carrapatoso--Mischler give a near-equilibrium Landau theorem with
Maxwell accommodation and quantitative convergence:

\url{https://arxiv.org/abs/2407.09031}
